\section{Kết luận và Khuyến nghị}

\subsection{Kết luận chung về đề án}

Đề án cá nhân kết thúc học phần Quản lý chất lượng giáo dục nghề nghiệp với chủ đề giải quyết bài toán thực tiễn tại Trường Cao đẳng nghề Kỹ thuật -- Công nghiệp A đã hoàn thành trọn vẹn các mục tiêu nghiên cứu và đáp ứng xuất sắc các chuẩn đầu ra được xác định trong đề cương môn học. 

Thông qua việc tiếp cận có hệ thống các văn bản quy phạm pháp luật và lý luận quản trị chất lượng hiện đại, đề án đã đạt được ba kết quả mang tính khoa học và thực tiễn sâu sắc:

Thứ nhất, đề án đã làm sáng tỏ căn cứ pháp lý và chẩn đoán toàn diện nguyên nhân khiến Trường Cao đẳng nghề Kỹ thuật -- Công nghiệp A không đạt tiêu chuẩn kiểm định cơ sở giáo dục nghề nghiệp năm 2026. Phân tích đối chiếu định lượng với Khoản 1 Điều 5 Thông tư số 14/2024/TT-BLĐTBXH đã khẳng định nhà trường vi phạm đồng thời cả hai điều kiện: vi phạm Điều kiện a do tổng điểm tự đánh giá chỉ đạt 79 trên 100 điểm, và vi phạm Điều kiện c do cả hai tiêu chí trụ cột phản ánh năng lực vận hành trực tiếp là Tiêu chí 2 về Hoạt động đào tạo và Tiêu chí 5 về Cơ sở vật chất, thiết bị đào tạo đều chỉ đạt 75.0\% điểm chuẩn, dưới ngưỡng sàn bắt buộc 80\%. Việc vận dụng các quan điểm Giá trị gia tăng, Phù hợp mục tiêu, Thỏa mãn khách hàng và quy tắc 4M đã bóc tách rõ nét những khoảng cách giữa điểm số hồ sơ tuân thủ và mức độ hài lòng thực tế của sinh viên cũng như doanh nghiệp, chỉ rõ điểm nghẽn nằm ở phương thức quản trị thụ động và sự thiếu đồng bộ của trang thiết bị xưởng thực hành.

Thứ hai, đề án đã thiết lập thành công bộ công cụ và khung danh mục minh chứng chuẩn hóa cho chương trình đào tạo trình độ cao đẳng nghề Điện công nghiệp bám sát Chương III Thông tư số 14/2024/TT-BLĐTBXH với quy mô 07 tiêu chí và 40 tiêu chuẩn. Việc tiến hành thẩm định thử minh chứng của năm tiêu chuẩn điển hình và biên soạn hoàn chỉnh Báo cáo tự đánh giá cho Tiêu chuẩn 1.3 về tích hợp năng lực số đã chứng minh tính khả thi trong việc kiểm chứng dữ liệu học thuật. Kết quả tính toán điểm tự đánh giá dự kiến đạt 82.5 trên 100 điểm, đáp ứng đầy đủ cả ba điều kiện tại Điều 11 của thông tư kiểm định, khẳng định chương trình đào tạo nghề Điện công nghiệp hoàn toàn đủ điều kiện pháp lý để tự tin bước vào đợt đánh giá ngoài độc lập, đồng thời cảnh báo các rủi ro cận biên để nhà trường có biện pháp phòng ngừa kịp thời.

Thứ ba, đề án đã đề xuất một mô hình thiết kế hệ thống đảm bảo chất lượng bên trong hoàn chỉnh cho nhà trường dựa trên 05 bước kinh điển của David Lim (2001), kết hợp với Kế hoạch hành động cải tiến chất lượng chu kỳ 01 năm theo chu trình Hoạch định -- Thực hiện -- Kiểm tra -- Cải tiến (PDCA). Kế hoạch này tập trung giải quyết tận gốc các điểm nghẽn của Tiêu chí 2 và Tiêu chí 5, thiết lập ma trận phân công trách nhiệm RACI minh bạch, định lượng các chỉ số hiệu năng cốt lõi và phân bổ ngân sách 3.5 tỷ đồng đầu tư thực chất cho xưởng thực hành, mở ra lộ trình vững chắc đưa toàn trường đạt chuẩn kiểm định cấp cơ sở và xây dựng văn hóa chất lượng tự giác, bền vững.

\subsection{Khuyến nghị thực thi}

Để các giải pháp trong đề án được chuyển hóa thành kết quả thực tế, tác giả đề xuất một số khuyến nghị trọng tâm đối với các chủ thể liên quan:

Đối với Ban Giám hiệu và Hội đồng trường, cần chuyển biến mạnh mẽ từ tư duy ``đối phó kiểm định theo đợt'' sang tư duy ``quản trị chất lượng toàn diện thường xuyên''. Lãnh đạo nhà trường cần cam kết bố trí nguồn lực tài chính ổn định hàng năm, dành tối thiểu từ 3\% đến 5\% tổng kinh phí chi thường xuyên cho các hoạt động đảm bảo chất lượng; duy trì nghiêm túc các phiên họp giao ban định kỳ hàng tháng của Hội đồng Đảm bảo chất lượng; đồng thời thực hiện phân quyền và trao quyền thực chất cho Phòng Đảm bảo chất lượng trong việc giám sát, cảnh báo sớm các sai lệch quy trình đào tạo.

Đối với Khoa Điện -- Điện tử, cần phát huy vai trò chủ động của đơn vị chuyên môn cơ sở trong việc rà soát chuẩn đầu ra và cấu trúc chương trình đào tạo. Đội ngũ nhà giáo trong khoa cần tích cực đổi mới phương pháp giảng dạy sang mô hình học tập dựa trên dự án, chủ động tự bồi dưỡng nâng cao năng lực ứng dụng phần mềm mô phỏng và kỹ năng số công nghiệp; tăng cường phối hợp chặt chẽ với các kỹ sư của doanh nghiệp đối tác trong việc đồng thiết kế bài giảng và tổ chức đánh giá kỹ năng thực hành cho sinh viên ngay tại hiện trường sản xuất.

Đối với Phòng Đảm bảo chất lượng và các phòng ban chức năng, cần khẩn trương số hóa toàn diện hệ thống quản lý minh chứng kiểm định, xây dựng cơ sở dữ liệu liên thông giữa Phòng Đào tạo, Phòng Quản trị Thiết bị, Trung tâm Hợp tác Doanh nghiệp và các khoa chuyên môn. Các đơn vị cần xóa bỏ hoàn toàn tâm lý e ngại cung cấp số liệu, quán triệt tinh thần minh bạch và coi dữ liệu phản hồi là công cụ khoa học để tự cải tiến công việc của chính đơn vị mình.

Đối với cơ quan quản lý nhà nước về giáo dục nghề nghiệp (Bộ Giáo dục và Đào tạo, Bộ Lao động -- Thương binh và Xã hội), kiến nghị tiếp tục hoàn thiện hành lang pháp lý đồng bộ trong giai đoạn thực thi Luật số 124/2025/QH15 (có hiệu lực từ ngày 01/01/2026 thay thế Luật Giáo dục nghề nghiệp số 74/2014/QH13) \parencite{luatgdnn2025}; ban hành cẩm nang hướng dẫn chi tiết về kỹ thuật số hóa, mã hóa và lưu trữ minh chứng điện tử theo Thông tư số 14/2024/TT-BLĐTBXH \parencite{bldtbxh2024}; tổ chức các chương trình tập huấn chuyên sâu về quản lý chất lượng và đảm bảo chất lượng cho đội ngũ cán bộ quản lý, nhà giáo tại các trường cao đẳng; đồng thời có chính sách ưu tiên đầu tư các gói dự án hiện đại hóa xưởng thực hành công nghệ cao cho các ngành nghề kỹ thuật trọng điểm quốc gia như nghề Điện công nghiệp.
